\section{Theoretical Analysis}
We provide a theoretical analysis showing that RCS is strictly superior to majority voting (self-consistency) under a Gaussian generative model on the embedding space.

\paragraph{Assumption of Gaussian Distribution.}
We model the embedding space \(\mathbb{R}^d\) (\(d \gg 1\)) where the \emph{embeddings} of the generated answers are given by
\begin{align}
\mathbf{u}_i = \mathbf{c}^* + \boldsymbol{\varepsilon}_i, \quad \boldsymbol{\varepsilon}_i \sim \mathcal{N}(\mathbf{0}, \sigma^2 \mathbf{I}), \quad i.i.d.,
\end{align}
with \(\mathbf{c}^*\) the unknown true semantic center (embedding of the highest-quality answer) and \(\boldsymbol{\varepsilon}_i\) capturing semantic noise such as paraphrasing or minor hallucinations. The discrete answers \(a_i\) themselves are produced by the LLM; the Gaussian assumption applies directly to their embeddings (the space relevant for RCS). Answer quality is inversely related to noise magnitude:
\[
Q_i = Q(a_i | x) \propto -\|\boldsymbol{\varepsilon}_i\|_2.
\]
The distribution \(P{=}\{p_i\}_{i=1}^N\) can be uniform, frequency-weighted, or probability-weighted (\(p_i \propto \exp(-\lambda \|\boldsymbol{\varepsilon}_i\|_2^2)\)), as defined in Table~\ref{tab:rds_variants}. The probability-weighted case arises naturally when the LLM is confidence-calibrated: its assigned probability decreases exponentially with semantic deviation, which under Gaussian noise yields exactly \(p_i \propto \exp(-\lambda \|\boldsymbol{\varepsilon}_i\|_2^2)\) with \(\lambda > 0\) controlling sharpness. We then have the following Propositions, detailed proofs are provided in Apppendix~\ref{app:proof_prop1} and~\ref{app:proof_prop2}.

% TODO: Move Props to Method, others in Appendix
\begin{prop}[\textbf{Asymptotic Superiority}]
Under the Gaussian generative model, for \emph{any} of the three distributions \(P\) (uniform, frequency-weighted, or probability-weighted), as \(N \to \infty\),
\begin{align}
\mathbf{c}(P) \xrightarrow{a.s.} \mathbf{c}^*,
\end{align}
by the strong law of large numbers. Consequently,
\begin{align}
\mathbb{P}(\text{RCS selects the true best answer}) \to 1,
\end{align}
while majority voting (which relies on exact string match) satisfies
\begin{align}
\mathbb{P}(\text{voting selects the true best answer}) \leq \frac{1}{M},
\end{align}
where \(M \geq 2\) is the number of distinct semantic clusters (typically \(M \gg 1\) in open-ended generation). Thus RCS is asymptotically strictly superior to majority voting whenever lexical variance (due to paraphrasing) exceeds semantic variance.
\end{prop}

\begin{prop}[\textbf{Finite-\(N\) Gap with 2-Cluster Toy Model}]
Consider the following two-cluster setting: (1) Cluster A (true high-quality): \(k \geq N/2\) samples centered at \(\mathbf{c}^* + \boldsymbol{\delta}\), \(\|\boldsymbol{\delta}\|_2\) small; and (2) Cluster B (spurious): \(N-k\) samples centered at \(\mathbf{c}^* + \boldsymbol{\Delta}\), \(\|\boldsymbol{\Delta}\|_2 \gg \|\boldsymbol{\delta}\|_2\).

Under the frequency-weighted distribution, the empirical center \(\mathbf{c}(P)\) is geometrically biased toward Cluster A when \(k \geq N/2\). Approximating each cluster's candidates at their expected locations, RCS selects Cluster A with probability
\begin{align}
\mathbb{P}(\text{RCS correct})
= \Phi\left( \frac{(2k-N)\bigl(\|\boldsymbol{\Delta}\|_2 - \|\boldsymbol{\delta}\|_2\bigr)}{2\sigma\sqrt{N}} \right),
\end{align}
where \(\Phi\) is the standard normal CDF. Under exact string matching with \(m \geq 2\) distinct surface forms per semantic cluster, a single random sample is correct with probability \(k/N\).

Therefore,
\begin{align}
\mathbb{P}(\text{RCS correct}) > \frac{k}{N}
\end{align}
whenever the semantic gap satisfies
\begin{align}
\|\boldsymbol{\Delta}\|_2 - \|\boldsymbol{\delta}\|_2 > c\frac{\sigma}{\sqrt{N}}, \quad c > 0.
\end{align}
This property holds for the uniform and probability-weighted variants. When \(k < N/2\), the center \(\mathbf{c}(P)\) is biased toward Cluster B; asymptotic correctness is still guaranteed by Proposition 1.
\end{prop}

% \subsection{When RCS Provably Outperforms Voting (Task Conditions)}
% RCS yields the largest theoretical and empirical gains under the following conditions:

% \begin{table}[h]
% \centering
% \begin{tabular}{lll}
% \toprule
% \textbf{Task type} & \textbf{Reason RCS \(\succ\) voting} & \textbf{Example} \\
% \midrule
% Open-ended / paraphrase-heavy & High lexical diversity, high semantic agreement & Explanations, creative QA \\
% Multi-step reasoning & Multiple semantically equivalent paths & GSM8K, arithmetic \\
% Hard expert QA & Minority high-quality answers near consensus & GPQA \\
% Low lexical overlap & Natural paraphrasing by LLM & SciQ (explanatory) \\
% Black-box / low \(N\) & No need for exact string match & All settings (\(N=10\)--\(20\)) \\
% \bottomrule
% \end{tabular}
% \caption{Tasks where the generative-model assumptions hold strongly and RCS gap is largest.}
% \label{tab:rcs_tasks}
% \end{table}